% Generated from chapter-03.md - do not edit.

\chapter{Better Thoughts In The Morning}

I don't wake up happy.

There. I've said it.

Perhaps there are people who open their eyes in the morning filled with enthusiasm for the new day.

If so, I would be interested in meeting them and learning their methods.

My own brain has a different morning routine.

Before I've even gotten out of bed, it likes to begin with something along the lines of:

\emph{Oh no.}

Then:

\emph{What time is it?}

Then perhaps:

\emph{I don't think I can do this.}

Do what?

Apparently \textbf{the day}.

I haven't even tried it yet.

\section{Don't Believe Everything You Think Before Breakfast}%
\index{morning}\index{before breakfast}\nobreak%

Morning is a particularly interesting time for me because my first thoughts often feel completely convincing.

I wake up and my body doesn't feel wonderful.

Something is stiff.

Something hurts.

I feel tired.

I wonder whether I slept enough.

I may have awakened from one of my peculiar dreams in which I am trying to go somewhere and can't get organized.

I'm packing for a trip but can't find my socks.

I've forgotten something important.

Everyone else seems to know what they're doing and I'm behind.

My dreams aren't usually terrifying. They're more like long demonstrations of my incompetence. Any dream interpretation I do leads to ``You are afraid.'' so I quit trying to figure out my mental state that way.

But I wake up carrying the feeling with me.

\emph{I'm failing.}

\emph{Something's wrong.}

\emph{This is going to be a bad day.}

And the remarkable thing is that \textbf{nothing has happened yet.}

I am lying safely in my own bed.

Nevertheless, my brain has prepared a complete review of the day and given it one star.

I have learned not to take that review too seriously.

\section{Seven O'Clock Me Is Not Ten O'Clock Me}

This has taken me a long time to learn.

How I feel when I first wake up is \textbf{not a reliable prediction of how I will feel later in the morning.}

That seems obvious when I write it.

It doesn't feel obvious at seven o'clock.

At seven o'clock, the thought is:

\emph{I feel awful. Therefore today will be awful.}

But I've lived through enough mornings now to have evidence.

I may feel miserable when I first wake up.

Then I lie there awhile.

I begin moving.

I take my pills.

I eat breakfast.

I listen to some music.

I do a little yoga or stretching.

I shower.

I get interested in something.

And somewhere along the way I discover:

\emph{Oh! I'm all right.}

Usually by nine o'clock I am downright cheerful.

What happened to the terrible day I predicted at seven?

It never arrived.

So now I have a better thought available:

\textbf{This is how I feel now. It is not necessarily how I will feel later.}

That little word \textbf{now} helps me enormously.

\emph{I'm tired now.}

\emph{I'm discouraged now.}

\emph{I'm worried now.}

That is quite different from:

\emph{I'm going to feel this way all day.}

I don't know that.

I've simply been awake for five minutes and gotten ahead of myself.

\section{The Morning Brain Is a Storyteller}%
\index{storytelling brain}\nobreak%

Chapter 2 was about the stories our minds create.

Morning gives me a perfect laboratory for watching mine at work.

I wake up with an uncomfortable feeling.

That's the fact.

Then the story begins.

\emph{Maybe I'm getting sick.}

\emph{Maybe I won't have enough energy.}

\emph{Maybe I won't get anything accomplished.}

\emph{Maybe I'll disappoint somebody.}

\emph{Maybe I won't be able to do what I planned.}

Within minutes, a vague sensation has acquired a plot.

This is when I need to remember:

\textbf{A feeling is real. The story I attach to it may not be.}

I don't need to argue with the feeling.

If I feel tired, I feel tired.

If my joints hurt, they hurt.

If I didn't sleep well, I didn't sleep well.

But I don't have to turn those facts into a prophecy.

Instead of:

\emph{This is going to be a terrible day.}

I can say:

\textbf{Let's see.}

That may be one of my favorite better thoughts.

\emph{Let's see.}

It doesn't demand optimism.

It doesn't deny anything.

It simply refuses to predict the ending before the day has begun.

\section{Nurture Yourself Awake}%
\index{nurture awake}\nobreak%

For years I thought I should wake up and \textbf{get going}.

Now I think perhaps I need to wake up and \textbf{get acquainted with myself first}.

How am I doing?

What does my body need?

What does my mind need?

What would make this beginning a little easier?

I don't want to bully myself out of bed.

I want to nurture myself awake.

So instead of:

\emph{Come on. Get up. You're wasting time.}

I might say:

\emph{You don't have to get up yet.}

\emph{Let yourself wake up.}

\emph{You're safe here.}

\emph{Rest a little longer.}

\emph{You will probably feel better after breakfast.}

That last one is especially useful because breakfast has rescued a surprising number of mornings.

Sometimes a better thought is extremely practical.

\section{Talk to Yourself as Someone You Love}%
\index{talk to yourself}\nobreak%

I have spent much of my life thinking about communication.

\index{Marriage Enrichment}\index{Parent Effectiveness Training}
I taught Parent Effectiveness Training. Jim and I spent decades in Marriage Enrichment. I learned about active listening and how much it matters to have another person really hear what we're feeling.

Eventually I realized:

\textbf{I am also in a relationship with myself.}

And sometimes I'm not particularly good company.

Imagine waking up beside someone who immediately said:

``You look terrible.''

``You didn't sleep enough.''

``You aren't going to get everything done.''

``You should have gone to bed earlier.''

``What's wrong with you?''

You might consider separate bedrooms.

Yet we can speak to ourselves that way without even noticing.

What if I use some of the communication skills I have practiced with other people on myself?

I can listen.

\emph{Oh, you're worried this morning.}

\emph{You're afraid you won't have enough energy.}

\emph{You're disappointed that you didn't sleep longer.}

\emph{You feel overwhelmed and you haven't even gotten out of bed.}

There is something soothing about simply being understood---even when the person doing the understanding is me.

And then perhaps:

\emph{What do you need?}

Not:

\emph{What's wrong with you?}

\textbf{What do you need?}

Those are very different questions.

\section{My Morning Delights}%
\index{morning delights}\nobreak%

Eventually I do get up.

And I have discovered that it helps enormously to have some things waiting for me that I actually like.

I think of these as my morning \textbf{Delights}. A better name than Chores.\index{morning delights}

Prayer.

A little yoga.

Exercises that help my back and shoulders.

Music.

A shower.

Taking care of my teeth.

Breakfast.

These are ordinary things.

That's precisely why they work.

They move me out of the imaginary day in my head and into the actual morning happening around me.

One thing.

Then another.

I don't have to decide whether I can manage the entire day.

I only have to do the next thing.

Sometimes the better thought is:

\textbf{I don't have to do Tuesday. I only have to get dressed.}

That's manageable.

\section{Give Yourself Something to Look Forward To}%
\index{look forward}\nobreak%

I've also discovered that tomorrow morning can sometimes be helped tonight.

If I go to sleep with nothing in particular to anticipate, my morning brain has an empty stage on which to perform.

And it knows a lot of tragedies.

So I like to give it something else.

Tomorrow I might take a walk.

I might see someone I love.

I might work on my book.

I might have a particularly good lunch.

I might read something I enjoy.

It doesn't have to be important.

Anticipation itself is pleasurable.

The night before, I can plant a little thought:

\textbf{There's something nice waiting for me tomorrow.}

Then perhaps my mind will find that thought in the morning before it finds something to worry about.

Perhaps.

My mind makes no guarantees.

\section{Start With a Memory}

One of my favorite morning practices grew directly out of what I learned after Jim died.

I deliberately remember something good.

Not necessarily from yesterday.

It can come from anywhere in my life.

A peak experience.\index{peak experience}

Something funny.

Something I accomplished.

A person I loved.

A moment when I was brave.

A time I helped somebody.

A beautiful place.

A completely insignificant little moment that still makes me smile.

I have lived a long time.

That means I have accumulated an enormous supply of memories.

Why leave them sitting in storage?

My brain is perfectly willing to retrieve an embarrassing thing I said forty years ago without being asked.

I think it is only fair that I ask it to retrieve some of the good material too.

So I choose one.

And I go into it.

What could I see?

Who was there?

What happened?

How did I feel?

I don't merely label it:

\emph{That was nice.}

I try to \textbf{visit it}.

For a little while, I let my mind live there.

\section{What Deserves My Attention?}%
\index{deserving attention}\nobreak%

There is a difference between ignoring a problem and deciding how much attention it deserves.

If I wake up with an important problem, perhaps I need to deal with it.

But do I need to deal with it at 6:17 in the morning while lying in bed?

Maybe not.

I can tell myself:

\textbf{I'll think about that after breakfast.}

Or:

\textbf{There is nothing I can do about that right now.}\index{nothing I can do}

Or:

\textbf{I've already decided what I'm going to do. I don't need another meeting about it inside my head.}

I particularly like that last one.

My mind can hold committee meetings that nobody requested.

Sometimes I need to adjourn them.

\section{When a Better Thought Isn't Available}

There are mornings when none of this works particularly well.

I want to say that clearly.

Sometimes I'm exhausted.

Sometimes my body doesn't feel good.

Sometimes something genuinely sad or frightening is happening.

Sometimes I ask, \emph{What's a better thought?} and my brain replies:

\emph{I have nothing.}

All right.

Then perhaps I don't need a better thought yet.

Perhaps I need breakfast.

Or rest.

Or prayer.

Or another person.

Or to sit quietly.

Or to cry.

Or to wait.

\textbf{Better thinking is not another assignment to fail.}

I don't want anyone reading this book to lie in bed feeling terrible and then add:

\emph{And now I'm failing at better thoughts.}

No.

Some mornings the better thought might simply be:

\textbf{I don't have to fix this right now.}

Or:

\textbf{I can be kind to myself while I wait for this feeling to pass.}

That counts.

\section{Borrow Someone Else's Brain}%
\index{someone else brain}\nobreak%

There are also mornings when I need another person.

I've learned that being heard can help loosen a thought just as it did all those years ago at that Marriage Enrichment retreat.

Sometimes I call a friend.

Sometimes I pray.

And in this rather remarkable period of history, sometimes I talk with artificial intelligence.

I can say what I'm thinking.

I can hear it reflected back.

I can ask for another way of looking at it.

Sometimes what I need most is not an answer.

I need help getting enough distance from my own thinking to see that there might be another thought available.

I still believe deeply in human connection.

But I also believe in using whatever good tools are available to help us care for our minds.

\section{The Sixteen-Minute Experiment}%
\index{sixteen minute experiment}\nobreak%

Recently I began experimenting with another idea.

I can set my phone or watch to remind me at intervals---sometimes every sixteen minutes.

The reminder isn't there to tell me I've done something wrong.

It asks me, in effect:

\textbf{Where is your mind?}

Perhaps I'm already perfectly happy.

Wonderful. Continue.

But perhaps I've wandered into worry.

Perhaps I'm rehearsing an argument.

Perhaps I'm criticizing myself.

Perhaps I've been staring at the fly on the masterpiece again.

The little sound interrupts me.

And I can decide:

\emph{What do I want to think about now?}

Love?

Goodness?

Someone I'm grateful for?

Something beautiful?

Jesus?

A memory?

Something I'm looking forward to?

The reminder doesn't create the thought.

It gives me something even more important:

\textbf{a moment of choice.}

I wish I had known to do something like this during the terrible months after Jim died.

But perhaps I had my own primitive version then.

A scrap of paper.

\textbf{What's a better thought?}

That little piece of paper was my reminder.

\section{We Don't Have to Choose the First Thought}

This may be the most important thing I have learned about mornings.

\textbf{I don't necessarily choose the first thought.}

It arrives.

Sometimes it has come out of a dream.

Sometimes it comes from my body.

Sometimes it is an old habit.

Sometimes I have no idea where it came from.

I don't have to be responsible for every thought that knocks on the door.

But I can become more interested in what happens \textbf{after} it arrives.

The first thought might be:

\emph{Oh no.}

The second can be:

\textbf{Good morning. Let's see how this goes.}

The first might be:

\emph{I feel terrible.}

The second:

\textbf{You often feel better later. Give yourself some time.}

The first:

\emph{I can't possibly do everything today.}

The second:

\textbf{You don't have to. What's the next thing?}

The first:

\emph{Something is wrong.}

The second:

\textbf{Maybe. Maybe not. Let's have breakfast before we decide.}

And sometimes the first thought is simply sad.

Then perhaps the second thought doesn't need to contradict it.

\emph{I miss him.}

And:

\textbf{Wasn't I lucky to have loved him?}

Both belong.

\section{Ten Percent Better}%
\index{ten percent better}\nobreak%

I don't need to wake up ecstatic.

That's asking quite a lot before coffee.

I don't even have to become happy.

Sometimes I ask myself a smaller question:

\textbf{What would make me feel ten percent better right now?}

Music?

Getting warm?

Opening the curtains?

A glass of water?

Prayer?

Breakfast?

A memory?

Getting out of bed?

Staying in bed another ten minutes?

Calling someone?

Ten percent counts.

And perhaps after that there is another ten percent.

This is not a dramatic transformation.

It is a direction.

\section{At the Threshold of the Day}%
\index{threshold of the day}\nobreak%

Every morning is a little beginning.

Before I know what the day will bring, there I am again---opening my eyes and meeting my mind.

I used to think the goal was to wake up happy.

Now I think the goal is something gentler.

\textbf{Wake up curious.}

Don't decide too soon.

Don't turn a feeling into a prophecy.

Don't mistake the first thought for the truth.

Listen to yourself.

Take care of your body.

Remember something wonderful.

Give yourself something to look forward to.

Borrow another brain when you need one.

And if the mind has already wandered down an old familiar track, don't scold it.

Notice.

Smile if you can.

And ask:

\textbf{What's a better thought---right now?}

Maybe there is one.

Maybe there isn't yet.

Either way, the day has only just begun.

\textbf{Let's be curious and see.}
